% Part: lambda-calculus
% Chapter: introduction
% Section: minimization

\documentclass[../../../include/open-logic-section]{subfiles}

\begin{document}

\olfileid{lam}{int}{min}
\olsection{\usetoken{S}{lambda definable} વિધેયો લઘુત્તમીકરણ હેઠળ બંધ છે}

\begin{lem}
ધારો કે $f(x,y)$ !!{lambda definable} છે. વિધેય~$g$ને
આ રીતે વ્યાખ્યાયિત કરો:
\[
g(x) \simeq \umin{y}{f(x,y)}.
\]
ત્યારે~$g$ પણ !!{lambda definable} છે.
\end{lem}

\begin{proof}
મુખ્ય વિચાર લગભગ આ છે. $x$ આપેલું હોય ત્યારે નિશ્ચિત-બિંદુ
લૅમ્બડા પદ~$Y$ વડે એવું વિધેય~$h_x(n)$ વ્યાખ્યાયિત કરીશું જે~$n$થી
શરૂ કરીને યોગ્ય~$y$ શોધે; પછી $g(x)$ એટલે માત્ર~$h_x(0)$.
વિધેય~$h_x$ને નિશ્ચિત-બિંદુ સમીકરણના ઉકેલ તરીકે લખી શકાય:
\[
h_x(n) \simeq
\begin{cases}
n & \text{જો $f(x,n) = 0$ હોય} \\
h_x(n+1) & \text{અન્યથા.}
\end{cases}
\]

હવે વિગતો જોઈએ. $f$ લૅમ્બડા-વ્યાખ્યેય હોવાથી કોઈ પદ~$F$
વડે તે !!{lambda defined} છે. યાદ કરો કે આપણી પાસે એવું લૅમ્બડા
પદ~$D$ પણ છે કે $D(M, N, \bar{0}) \red M$ અને $D(M, N, \bar{1})
\red N$. હાલ પૂરતું~$x$ નિશ્ચિત રાખીએ. $h_x$ને
!!{lambda define} કરવા માટે~$x$ પર આધાર રાખતું એવું પદ~$H$
શોધવું છે જે નીચેની શરત પૂરી કરે:
\[
H(\num{n}) \equiv D(\num{n}, H(S(\num{n})), F(x, \num{n})).
\]
નિશ્ચિત-બિંદુ પદ~$Y$ વડે આ કરી શકાય છે. પહેલાં~$U$ને નીચેનું
પદ માનો:
\[
\lambd[h][\lambd[z][D(z,(h(Sz)),F(x,z))]],
\]
અને પછી~$H$ને પદ~$YU$ માનો. ધ્યાન આપો કે~$H$માં એકમાત્ર મુક્ત
ચલ~$x$ છે. હવે~$H$ ઉપરનું સમીકરણ પૂરે છે તે બતાવીએ.
\sourcecorrection{OLLAM-009}{લેમ્માની ધારણા માત્ર એટલી છે કે~$f$
લૅમ્બડા-વ્યાખ્યેય છે. સ્થિર અંગ્રેજી મૂળ અહીં અચાનક~$f$ને આદિમ
પુનરાવર્તી કહે છે, જે અનાવશ્યક મજબૂત ધારણા છે. આપેલી ધારણાથી
જ~$F$ મળે છે એમ ગુજરાતી પાઠમાં લખ્યું છે.}

$Y$ની વ્યાખ્યા મુજબ,
\[
H = YU \equiv U(YU) = U(H).
\]
ખાસ કરીને, દરેક પ્રાકૃતિક સંખ્યા~$n$ માટે,
\begin{align*}
H(\num{n}) & \equiv U(H, \num{n}) \\
& \red D(\num{n}, H(S(\num{n})), F(x, \num{n})),
\end{align*}
જે જોઈએ હતું તે જ. છેલ્લી પંક્તિમાં~$x$ની જગ્યાએ અંક~$\num{m}$
મૂકો. જો $F(\num{m}, \num{n})$નું લઘુકરણ~$\num{0}$ સુધી થાય,
તો આખી અભિવ્યક્તિ~$\num{n}$ સુધી ઘટે છે; અને જો
$F(\num{m}, \num{n})$નું લઘુકરણ કોઈ બીજા અંક સુધી થાય,
તો અભિવ્યક્તિ~$H(S(\num{n}))$ સુધી ઘટે છે.
આ બીજો કિસ્સો અંક એક પૂરતો મર્યાદિત નથી: જોડી-પસંદગીમાં શૂન્યથી
મોટા ચર્ચ અંક પર અચલ વિધેય વારંવાર લાગુ પડે છે અને બીજો ઘટક જ આપે છે.

સાબિતી પૂરી કરવા~$G$ને $\lambd[x][H(\num 0)]$ માનો.
ત્યારે~$G$, $g$ને !!{lambda define} કરે છે. બીજા શબ્દોમાં,
દરેક~$m$ માટે, જો~$g(m)$ વ્યાખ્યાયિત હોય તો
$G(\num m)$નું લઘુકરણ~$\overline {g(m)}$ સુધી થાય છે;
નહીં તો તેનું કોઈ સામાન્ય સ્વરૂપ હોતું નથી.
\end{proof}

\end{document}
